% =============================================================
%  CS102-F26 Digital Design and Computer Architecture I
%  Lecture 02 -- Numeral Systems (Unit 2)
%  Follows Section 2.1 of A Student's Computer Architecture Book I.
%
%  Build the slides and the print handout:  npm run materials -- lecture-02
% =============================================================
\makeatletter
\def\input@path{{templates/latex/}{../}}
\makeatother

\documentclass[10pt]{wosm-slides}
\usepackage{cs102-lectures}

\title[Lecture 02 \textperiodcentered\ Numeral Systems]{Numeral Systems}
\subtitle{Lecture 02: additive and positional systems; binary, octal and hexadecimal}

% ---------- Helpers for this lecture --------------------------
% A numeral in a given base, digits upright: \base{2F}{16}.
\newcommand*\base[2]{\ensuremath{\mathrm{#1}_{#2}}}

% Steps of a conversion, one row each.
\newenvironment{steps}[1]
  {\begin{tabular}{@{}#1@{}}\toprule}
  {\bottomrule\end{tabular}}

% =============================================================
\begin{document}

\maketitle

\begin{frame}{Today}
  \tableofcontents
\end{frame}

% =============================================================
\section{Ways of counting}

\begin{frame}{Numbers and numerals}
  Counting is an unavoidable part of human life, and people have invented many
  different ways of writing counts down.

  \begin{definition}[Numeral system]
    A \emph{numeral system} is a way of writing numbers with symbols. A
    \emph{numeral} is one number written in that way.
  \end{definition}

  \medskip
  One number, four numerals:
  \begin{center}
    \begin{tabular}{@{}cccc@{}}
      $|\,|\,|\,|\,|\,|$ & VI & 6 & \base{110}{2} \\[2pt]
      \wosmlabel{tally} & \wosmlabel{Roman} & \wosmlabel{decimal} & \wosmlabel{binary} \\
    \end{tabular}
  \end{center}

  \medskip
  There are two big families: \alert{additive} and \alert{positional} systems.
\end{frame}

\begin{frame}{Additive systems}
  In an additive system, each symbol has a fixed value, and the number is the sum
  of the symbols. Where a symbol stands does not change its magnitude.

  \medskip
  \begin{columns}[T, totalwidth=\textwidth]
    \column{0.4\textwidth}
      \textbf{Roman numerals}\par\smallskip
      \footnotesize
      \begin{tabular}{@{}*{7}{c}@{}}
        \toprule
        I & V & X & L & C & D & M \\
        \midrule
        1 & 5 & 10 & 50 & 100 & 500 & 1000 \\
        \bottomrule
      \end{tabular}
    \column{0.56\textwidth}
      \begin{itemize}
        \item $\mathrm{VI} = 5 + 1 = 6$
        \item $\mathrm{IV} = -1 + 5 = 4$: a smaller symbol before a larger one is
          subtracted, but its magnitude is still 1
        \item $\mathrm{MMXXVI} = 1000 + 1000 + 10 + 10 + 5 + 1 = 2026$
      \end{itemize}
  \end{columns}

  \bigskip
  Tally marks are additive too: every mark is worth 1, wherever it is.
\end{frame}

\begin{frame}{The trouble with additive systems}
  \begin{itemize}
    \item There is no symbol for zero.
    \item Larger numbers need new symbols, or very long numerals:
      $3888 = \mathrm{MMMDCCCLXXXVIII}$ takes 15 symbols.
    \item Written arithmetic is awkward. Try multiplying XLVII by XXXVIII on paper.
  \end{itemize}

  \bigskip
  \begin{remark}
    Positional systems are far more useful in almost every area of mathematics and
    in digital logic design. From now on we work only with them.
  \end{remark}
\end{frame}

\begin{frame}{Positional systems}
  In a positional system, the \alert{place} of a digit decides its magnitude.

  \medskip
  \begin{center}
    \begin{tabular}{@{}r@{\quad}c@{\quad}l@{}}
      123 & $=$ & $100 + 20 + 3$ \\
      321 & $=$ & $300 + 20 + 1$ \\
    \end{tabular}
  \end{center}
  The same three digits, three different numbers.
  \begin{itemize}
    \item Zero holds a place: 105 and 15 differ only by the 0.
    \item Our decimal digits are positional with ten digits. Other choices are just
      as possible.
  \end{itemize}

  \begin{example}[Base 60, still in use]
    The Babylonians counted in base 60, and we still do for time:
    1\,h 15\,min 30\,s is $1 \cdot 60^2 + 15 \cdot 60 + 30 = 4530$ seconds.
  \end{example}
\end{frame}

% =============================================================
\section{The radix}

\begin{frame}{Radix and digits}
  \begin{definition}[Radix]
    The \emph{radix}, or \emph{base}, $b$ of a positional system is the number of
    elementary symbols, the \emph{digits}, that all its numerals are made of. The
    digits stand for $0, 1, \dots, b - 1$.
  \end{definition}

  Each digit is multiplied by the power of $b$ for its position, counted from 0 on
  the right, and the results are added:
  \[
    \overline{d_n d_{n-1} \dots d_0}_{\,b}
      = d_n \cdot b^n + d_{n-1} \cdot b^{n-1} + \dots + d_0 \cdot b^0 .
  \]

  \begin{example}[Decimal, $b = 10$]
    $2026 = 2 \cdot 10^3 + 0 \cdot 10^2 + 2 \cdot 10^1 + 6 \cdot 10^0$
  \end{example}
\end{frame}

\begin{frame}{Binary, octal and hexadecimal}
  \begin{center}
    \begin{tabular}{@{}llll@{}}
      \toprule
      System & Radix & Digits & Example \\
      \midrule
      binary & 2 & 0, 1 & $\base{101101}{2} = 45$ \\
      octal & 8 & 0, 1, \dots, 7 & $\base{55}{8} = 45$ \\
      decimal & 10 & 0, 1, \dots, 9 & $\base{45}{10} = 45$ \\
      hexadecimal & 16 & 0, 1, \dots, 9, A, B, C, D, E, F & $\base{2D}{16} = 45$ \\
      \bottomrule
    \end{tabular}
  \end{center}

  \medskip
  \begin{itemize}
    \item Hexadecimal needs sixteen digits, for the values 0 to 15. After 9 come
      the letters: A $= 10$, B $= 11$, \dots, F $= 15$.
    \item The value formula is the same for every radix:
      $\overline{b_n \dots b_0}_{\,2} = b_n \cdot 2^n + \dots + b_0$, and so on.
  \end{itemize}
\end{frame}

\begin{frame}{Counting in four bases}
  \small
  \begin{columns}[T, totalwidth=\textwidth]
    \column{0.47\textwidth}
      \centering
      \begin{tabular}{@{}rrrr@{}}
        \toprule
        dec & bin & oct & hex \\
        \midrule
        0 & 0000 & 0 & 0 \\
        1 & 0001 & 1 & 1 \\
        2 & 0010 & 2 & 2 \\
        3 & 0011 & 3 & 3 \\
        4 & 0100 & 4 & 4 \\
        5 & 0101 & 5 & 5 \\
        6 & 0110 & 6 & 6 \\
        7 & 0111 & 7 & 7 \\
        \bottomrule
      \end{tabular}
    \column{0.47\textwidth}
      \centering
      \begin{tabular}{@{}rrrr@{}}
        \toprule
        dec & bin & oct & hex \\
        \midrule
        8 & 1000 & 10 & 8 \\
        9 & 1001 & 11 & 9 \\
        10 & 1010 & 12 & A \\
        11 & 1011 & 13 & B \\
        12 & 1100 & 14 & C \\
        13 & 1101 & 15 & D \\
        14 & 1110 & 16 & E \\
        15 & 1111 & 17 & F \\
        \bottomrule
      \end{tabular}
  \end{columns}

  \medskip
  \normalsize
  In every base, counting works like an odometer: when a digit passes its largest
  value, it goes back to 0 and the next digit goes up by one ($\base{1111}{2} + 1 =
  \base{10000}{2} = 16$).
\end{frame}

\begin{frame}{Why binary?}
  A switch is on or off: two states, so two digits (Lecture 01).

  \begin{itemize}
    \item One binary digit is a \alert{bit}. Eight bits are a \alert{byte}; four bits
      are a \emph{nibble}, exactly one hexadecimal digit.
    \item $n$ bits have $2^n$ states, so they can stand for the numbers $0$ to
      $2^n - 1$.
  \end{itemize}

  \medskip
  \begin{center}
    \small
    \begin{tabular}{@{}l*{11}{r}@{}}
      \toprule
      $n$ & 0 & 1 & 2 & 3 & 4 & 5 & 6 & 7 & 8 & 9 & 10 \\
      $2^n$ & 1 & 2 & 4 & 8 & 16 & 32 & 64 & 128 & 256 & 512 & 1024 \\
      \bottomrule
    \end{tabular}
  \end{center}

  \medskip
  How many bits do you need for the numbers 0 to 999?
  \uncover<2->{$2^9 = 512$ is too few and $2^{10} = 1024$ is enough: 10 bits.}
\end{frame}

\begin{frame}{Writing the base}
  The digits alone do not say which base is meant: is 10 ten, two, eight or
  sixteen?

  \medskip
  \begin{center}
    \begin{tabular}{@{}lllll@{}}
      \toprule
      Where & Binary & Octal & Decimal & Hexadecimal \\
      \midrule
      these slides & \base{1011}{2} & \base{13}{8} & 11 & \base{B}{16} \\
      many programming languages & \texttt{0b1011} & \texttt{0o13} & \texttt{11}
        & \texttt{0xB} \\
      Verilog (4 bits) & \texttt{4'b1011} & \texttt{4'o13} & \texttt{4'd11}
        & \texttt{4'hB} \\
      \bottomrule
    \end{tabular}
  \end{center}

  \medskip
  The book writes $\overline{d_n \dots d_0}_{\,b}$, with a bar over the digits.
  Numbers without a base are decimal.
\end{frame}

% =============================================================
\section{Converting between bases}

\begin{frame}{To decimal: expand}
  Multiply each digit by its power of the base and add.

  \medskip
  \begin{align*}
    \base{101101}{2} &= 1 \cdot 32 + 0 \cdot 16 + 1 \cdot 8 + 1 \cdot 4
      + 0 \cdot 2 + 1 \cdot 1 = 45 \\
    \base{57}{8} &= 5 \cdot 8 + 7 = 47 \\
    \base{2F}{16} &= 2 \cdot 16 + 15 = 47
  \end{align*}

  \medskip
  \begin{remark}[A shortcut for binary]
    Every binary digit is 0 or 1, so just add up the powers of 2 where the bits are
    1: $\base{101101}{2} = 32 + 8 + 4 + 1 = 45$.
  \end{remark}
\end{frame}

\begin{frame}{To decimal: Horner's rule}
  Read the digits from the left. Keep one number $v$, starting at 0, and for each
  digit $d$: \quad $v \leftarrow v \cdot b + d$.

  \medskip
  \begin{columns}[T, totalwidth=\textwidth]
    \column{0.4\textwidth}
      \small
      \begin{steps}{lrl}
        Digit & $v$ & \\
        \midrule
        1 & 1 & $= 0 \cdot 2 + 1$ \\
        0 & 2 & $= 1 \cdot 2 + 0$ \\
        1 & 5 & $= 2 \cdot 2 + 1$ \\
        1 & 11 & $= 5 \cdot 2 + 1$ \\
        0 & 22 & $= 11 \cdot 2 + 0$ \\
        1 & 45 & $= 22 \cdot 2 + 1$ \\
      \end{steps}
    \column{0.54\textwidth}
      \begin{itemize}
        \item Converting \base{101101}{2}: no powers to compute, one
          multiplication and one addition per digit.
        \item A computation in the sense of Lecture 01: the state is $v$, the input
          is the next digit.
      \end{itemize}
  \end{columns}

  \medskip
  It works because the powers can be factored out:\par\smallskip
  \centering
  $\base{101101}{2} = ((((1 \cdot 2 + 0) \cdot 2 + 1) \cdot 2 + 1) \cdot 2 + 0)
   \cdot 2 + 1$\par
\end{frame}

\begin{frame}{From decimal: repeated division}
  Divide by the base and write down the remainder; repeat with the quotient until
  it is 0. The remainders, read from the \alert{last to the first}, are the digits.

  \medskip
  \begin{columns}[T, totalwidth=\textwidth]
    \column{0.4\textwidth}
      \small
      \begin{steps}{rrr}
        $N$ & $N \div 2$ & remainder \\
        \midrule
        45 & 22 & 1 \\
        22 & 11 & 0 \\
        11 & 5 & 1 \\
        5 & 2 & 1 \\
        2 & 1 & 0 \\
        1 & 0 & 1 \\
      \end{steps}
    \column{0.54\textwidth}
      \begin{itemize}
        \item Read upwards: $45 = \base{101101}{2}$.
        \item Why it works: $N = q \cdot b + d_0$, so the remainder is the last
          digit $d_0$, and the quotient $q$ has the remaining digits
          $d_n \dots d_1$.
        \item The same method works for any base.
      \end{itemize}
  \end{columns}
\end{frame}

\begin{frame}{Example: 2026 in hexadecimal}
  \begin{columns}[T, totalwidth=\textwidth]
    \column{0.44\textwidth}
      \small
      \begin{steps}{rrrl}
        $N$ & $N \div 16$ & remainder & \\
        \midrule
        2026 & 126 & 10 & A \\
        126 & 7 & 14 & E \\
        7 & 0 & 7 & 7 \\
      \end{steps}
    \column{0.5\textwidth}
      Read upwards: $2026 = \base{7EA}{16}$.

      \medskip
      Check by expanding:
      \[
        7 \cdot 256 + 14 \cdot 16 + 10 = 1792 + 224 + 10 = 2026 .
      \]
  \end{columns}

  \bigskip
  \begin{warning}[Common mistake]
    Reading the remainders from the top gives \base{AE7}{16}, which is 2791. The
    first remainder is the \emph{last} digit.
  \end{warning}
\end{frame}

\begin{frame}{Binary to octal and hexadecimal: group the bits}
  Since $8 = 2^3$ and $16 = 2^4$, every octal digit is exactly 3 bits and every
  hexadecimal digit exactly 4 bits. Group the bits from the \alert{right}.

  \medskip
  Take \base{1011010111}{2}:
  \begin{center}
    \begin{tabular}{@{}lccccc@{}}
      \toprule
      groups of 4 & & 0010 & 1101 & 0111 & \\
      hexadecimal & & 2 & D & 7 & $= \base{2D7}{16}$ \\
      \midrule
      groups of 3 & 001 & 011 & 010 & 111 & \\
      octal & 1 & 3 & 2 & 7 & $= \base{1327}{8}$ \\
      \bottomrule
    \end{tabular}
  \end{center}

  \medskip
  All three are 727. Pad the leftmost group with zeros if it is short.
\end{frame}

\begin{frame}{And back again}
  \begin{itemize}
    \item Hexadecimal to binary: write each digit as 4 bits.\par
      $\base{C3}{16} = \texttt{1100}\;\texttt{0011} = \base{11000011}{2} = 195$
    \item Octal to binary: write each digit as 3 bits.\par
      $\base{57}{8} = \texttt{101}\;\texttt{111} = \base{101111}{2} = 47$
    \item Octal to hexadecimal, or back: go through binary.\par
      $\base{1327}{8} = \base{1011010111}{2} = \base{2D7}{16}$
  \end{itemize}

  \bigskip
  \begin{remark}
    This is why hexadecimal is everywhere in computing: it is a short way of writing
    bits. One byte is always two hexadecimal digits, from \base{00}{16} to
    \base{FF}{16}.
  \end{remark}
\end{frame}

\begin{frame}{Try it}
  \begin{enumerate}
    \item Write 100 in binary, octal and hexadecimal.\par
      \uncover<2->{\wosmlabel{Answer:} $100 = 64 + 32 + 4 = \base{1100100}{2}$;
      $\texttt{1}\;\texttt{100}\;\texttt{100} = \base{144}{8}$;
      $\texttt{0110}\;\texttt{0100} = \base{64}{16}$.}
    \item What is \base{FF}{16} in binary and in decimal?\par
      \uncover<3->{\wosmlabel{Answer:} $\base{11111111}{2} = 2^8 - 1 = 255$, the
      largest number in one byte.}
    \item Which is larger: \base{1000000}{2} or \base{77}{8}?\par
      \uncover<4->{\wosmlabel{Answer:} $\base{1000000}{2} = 64$ and
      $\base{77}{8} = 63$, so \base{1000000}{2}, by one.}
  \end{enumerate}
\end{frame}

% =============================================================
\section{Numbers in the DX8}

\begin{frame}{Numbers in the DX8}
  The computer we build in this module works with bits everywhere:
  \begin{itemize}
    \item it is an 8-bit computer: one byte holds a number from 0 to 255;
    \item its 8 registers are selected by 3-bit addresses, \base{000}{2} to
      \base{111}{2}, because $2^3 = 8$;
    \item its RAM has $2^8 = 256$ bytes, so a memory address is one byte,
      \base{00}{16} to \base{FF}{16};
    \item an instruction is 16 bits: four hexadecimal digits.
  \end{itemize}

  \bigskip
  \begin{remark}[Coming back]
    So far every number is 0 or positive. Negative numbers in binary, and how to add
    them, come with two's complement \unitref{16}.
  \end{remark}
\end{frame}

% =============================================================
\section{Summary}

\begin{frame}{Summary}
  \begin{itemize}
    \item In an additive system each symbol has a fixed value; in a positional
      system the place of a digit decides its magnitude.
    \item With radix $b$: digits $0, \dots, b - 1$ and
      $\overline{d_n \dots d_0}_{\,b} = \sum_i d_i \cdot b^i$.
    \item To decimal: expand, or use Horner's rule $v \leftarrow v \cdot b + d$.
      From decimal: divide repeatedly, read the remainders from the last.
    \item Binary to octal or hexadecimal: group the bits by 3 or by 4, from the
      right.
    \item $n$ bits give $2^n$ numbers, from $0$ to $2^n - 1$.
  \end{itemize}
\end{frame}

\begin{frame}{Before the next lecture}
  \begin{itemize}
    \item Read Section~2.1 of the book, \emph{Numeral Systems}.
    \item Practise: write your year of birth in binary, octal and hexadecimal, and
      check each answer by converting it back.
    \item Count from 0 to 31 in binary. Which bit changes at every step, and which
      only once?
  \end{itemize}

  \begin{block}{Next lecture \unitref{3}}
    Boolean algebra: an algebra with only the values 0 and 1, and its laws and
    theorems.
  \end{block}
\end{frame}

\end{document}
